% !TeX root = ../../tesis.tex
% §2.1.5 — síntesis nueva (D36-4) + 2-5-SIG-Datos-Abiertos.tex l. 82-102 (D36-6) (#36)

% -----------------------------------------------------------------------
\subsection{Síntesis: Tendencias, Debates y Vacíos}
\label{subsec:ea-sintesis}
% -----------------------------------------------------------------------

La tendencia más clara es tratar la red de transporte como un grafo construido
con datos \gtfs{}. \citet{fortin2016innovative} cargan los \termino{feeds} en
una base de grafos para calcular caminos mínimos e indicadores;
\citet{kujala2018} normalizan con el mismo estándar las redes de 25 ciudades, y
\citet{ugur2025} miden la conectividad estación por estación. La segunda
tendencia es el \termino{software} propio y abierto. \citet{moreno2023}
defienden el desarrollo modular en Python para que una institución acumule
capacidad de análisis sin pagar licencias, y la cartografía colaborativa ya
cubre buena parte de la red vial de las ciudades latinoamericanas grandes
\citep{barrington2017world}.

Hay dos debates abiertos. El primero es cómo resumir varias dimensiones de una
red. Un índice compuesto da una sola cifra, pero un valor alto en una dimensión
puede tapar uno bajo en otra \citep{hossain2024}, y en redes multimodales la
concentración de flujos en pocos nodos se pierde dentro del índice
\citep{deng2023}. Por eso \citet{kujala2018} prefirieron un perfil por ciudad,
y esta tesis toma el mismo partido con la Clasificación Garibelt
(§\ref{subsubsec:garibelt-compensabilidad}). El segundo debate enfrenta la
demanda con la estructura. El Programa Maestro del Metro de 1985 midió con
detalle la demanda; con ella, los corredores del Periférico no alcanzaban la
captación que se exigía, y la evaluación no podía registrar lo que un anillo
aporta a la estructura de la red (§\ref{subsubsec:anillo-pm-2010}). Esta tesis
mide la estructura y deja fuera la demanda, de modo que ninguno de los dos
enfoques basta solo.

Quedan tres vacíos. El primero: de los anillos de otras metrópolis se conoce el
propósito, pero no hay una medición de la red (§\ref{subsec:ea-anillo-radial}).
El segundo está en los datos. \citet{haklay2010good} y
\citet{barrington2017world} miden la completitud de la red vial, no la de las
paradas de transporte, y en la \zmvm{} los sistemas que operan solo en el
Estado de México no tienen \gtfs{} oficial (§\ref{subsec:limitaciones-datos}).
El tercero es de herramientas. En México el \termino{software} propio de
análisis apenas empieza; el del Instituto Mexicano del Transporte se presenta
como primera etapa y trata la asignación de tráfico \citep{moreno2023}. Es la
brecha que describe \citet{gurstein2011open}: los datos abiertos sirven poco si
no hay herramientas para quien no programa.

\label{subsec:datos-herramienta-replicable}
Esta investigación responde a esa brecha con una cadena de tres motores de
código abierto. \texttt{Apímetro} reúne y normaliza los \termino{feeds}
\gtfs{}, \texttt{VFTModel} calcula los indicadores y
\texttt{Transport-gis-zmvm-mjg} los lleva a mapas que un urbanista puede leer sin
saber programar. Ninguno depende de licencias comerciales, y la cadena se puede
repetir en cualquier área metropolitana con \termino{feeds} \gtfs{} y
cartografía \gis{} de base, algo que, según \citet{barrington2017world}, cumple
la mayoría de las ciudades latinoamericanas de escala comparable a la \zmvm{}.

El argumento de replicabilidad no es secundario, sino central al objetivo
general de la investigación: demostrar que un planificador urbano con acceso a
datos abiertos puede, mediante esta metodología, obtener un diagnóstico
topológico riguroso de su red de transporte y evaluar escenarios alternativos de
intervención sin depender de \termino{software} propietario ni de equipos técnicos
especializados. Los referentes de calidad de datos
\citep{haklay2010good, barrington2017world} y los principios de equidad en el
acceso al conocimiento \citep{gurstein2011open, neis2012street} constituyen, en
este sentido, el fundamento teórico que articula la dimensión metodológica de la
tesis con su dimensión urbanística.
